% GENERATED FILE. Edit canonical lessons, then run npm run generate:books.
% canonical-source-hash: fnv1a64:3aa3f98a056ddf84
% canonical-lessons: RU-C456-obmennik, RU-C456-valyuta, RU-C456-garantiya, RU-C456-brakovannyy, RU-C456-ekonomit

\chapter{Currency and warranties}
\label{ch:ru-currency-and-warranties}
\hlchaptermodality{456}

\begin{chapteropening}
\textbf{By the end of this chapter:} \emph{I can deal with money and guarantees: a bureau de change, foreign currency, a warranty, faulty goods and saving money.}

Complete the last lesson of chapter 456: I can deal with money and guarantees: a bureau de change, foreign currency, a warranty, faulty goods and saving money.
\end{chapteropening}

\section[\texorpdfstring{\ru{обменник} (obménnik)}{обменник (obménnik)}]{\ru{обменник} (obménnik) — a bureau de change}

\label{lesson:RU-C456-obmennik}



\begin{quote}
Before the new one: say the Russian for to break, then the Russian for to shortchange.
\end{quote}

\subsection*{You'll want to know: \ru{обменник}}
\textbf{\ru{обменник}} (\emph{obménnik}) — \textquotedblleft{}a bureau de change\textquotedblright{}.

\begin{quote}
\ru{В} \ru{аэропорту} \ru{есть} \ru{обменник}.
\end{quote}

\begin{quote}
\emph{V aeroportú yest obménnik.}
\end{quote}

\begin{quote}
There is a bureau de change at the airport.
\end{quote}

Colloquial; officially \emph{\ru{пункт} \ru{обмена}}.

\begin{cousinweb}
\emph{\ru{обменник}} is built on \emph{\ru{обмен}} 'an exchange' with \emph{-\ru{ник}}, from \emph{\ru{менять}} 'to change', on \emph{\ru{мена}} 'an exchange', which is inherited from Proto-Slavic \emph{měna} (EtymDB).
\end{cousinweb}

\begin{grammarlens}[title={What you've built}]
One more word for telling, explaining and describing at B1.
\end{grammarlens}

\subsection*{Your turn}
\begin{itemize}
\raggedright
  \item \emph{Say it:} \emph{obménnik}
  \item \emph{Say it:} \emph{obménnik}, once more
  \item \emph{Say it:} a sentence of your own with \emph{obménnik}
  \item \emph{Recall:} say the Russian for to break, then the Russian for to shortchange, then say \emph{obménnik} again
\end{itemize}

\begin{tcolorbox}[breakable,colback=teal!4,colframe=teal!35!black,title={Before you move on}]
What does \ru{обменник} mean? (\textquotedblleft{}A bureau de change\textquotedblright{}.) Say it once more.
\end{tcolorbox}
\section[\texorpdfstring{\ru{валюта} (valyúta)}{валюта (valyúta)}]{\ru{валюта} (valyúta) — a foreign currency, a currency}

\label{lesson:RU-C456-valyuta}



\begin{quote}
Before the new one: say the Russian for to shortchange, then the Russian for a bureau de change.
\end{quote}

\subsection*{You'll want to know: \ru{валюта}}
\textbf{\ru{валюта}} (\emph{valyúta}) — \textquotedblleft{}a foreign currency, a currency\textquotedblright{}.

\begin{quote}
\ru{Мы} \ru{поменяли} \ru{валюту} \ru{в} \ru{банке}.
\end{quote}

\begin{quote}
\emph{My pomenyáli valyútu v bánke.}
\end{quote}

\begin{quote}
We changed currency at the bank.
\end{quote}

Feminine. In Russia \emph{\ru{валюта}} usually means foreign money.

\begin{cousinweb}
\emph{\ru{валюта}} was borrowed from Italian \emph{valuta} 'currency' (EtymDB), which goes back to Latin \emph{valēre} 'to be worth'.
\end{cousinweb}

\begin{grammarlens}[title={What you've built}]
One more word for telling, explaining and describing at B1.
\end{grammarlens}

\subsection*{Your turn}
\begin{itemize}
\raggedright
  \item \emph{Say it:} \emph{valyúta}
  \item \emph{Say it:} \emph{valyúta}, once more
  \item \emph{Say it:} a sentence of your own with \emph{valyúta}
  \item \emph{Recall:} say the Russian for to shortchange, then the Russian for a bureau de change, then say \emph{valyúta} again
\end{itemize}

\begin{tcolorbox}[breakable,colback=teal!4,colframe=teal!35!black,title={Before you move on}]
What does \ru{валюта} mean? (\textquotedblleft{}A foreign currency, a currency\textquotedblright{}.) Say it once more.
\end{tcolorbox}
\section[\texorpdfstring{\ru{гарантия} (garántiya)}{гарантия (garántiya)}]{\ru{гарантия} (garántiya) — a warranty, a guarantee}

\label{lesson:RU-C456-garantiya}



\begin{quote}
Before the new one: say the Russian for a bureau de change, then the Russian for a foreign currency.
\end{quote}

\subsection*{You'll want to know: \ru{гарантия}}
\textbf{\ru{гарантия}} (\emph{garántiya}) — \textquotedblleft{}a warranty, a guarantee\textquotedblright{}.

\begin{quote}
\ru{На} \ru{этот} \ru{телефон} \ru{гарантия} \ru{два} \ru{года}.
\end{quote}

\begin{quote}
\emph{Na étot telefón garántiya dva góda.}
\end{quote}

\begin{quote}
This phone has a two-year warranty.
\end{quote}

Feminine. \emph{\ru{По} \ru{гарантии}} — under warranty.

\begin{cousinweb}
\emph{\ru{гарантия}} is usually traced to French \emph{garantie} 'a guarantee' (EtymDB lists the French word), from Old French \emph{garantir} 'to vouch for', a word of Germanic origin.
\end{cousinweb}

\begin{grammarlens}[title={What you've built}]
One more word for telling, explaining and describing at B1.
\end{grammarlens}

\subsection*{Your turn}
\begin{itemize}
\raggedright
  \item \emph{Say it:} \emph{garántiya}
  \item \emph{Say it:} \emph{garántiya}, once more
  \item \emph{Say it:} a sentence of your own with \emph{garántiya}
  \item \emph{Recall:} say the Russian for a bureau de change, then the Russian for a foreign currency, then say \emph{garántiya} again
\end{itemize}

\begin{tcolorbox}[breakable,colback=teal!4,colframe=teal!35!black,title={Before you move on}]
What does \ru{гарантия} mean? (\textquotedblleft{}A warranty, a guarantee\textquotedblright{}.) Say it once more.
\end{tcolorbox}
\section[\texorpdfstring{\ru{бракованный} (brakóvannyy)}{бракованный (brakóvannyy)}]{\ru{бракованный} (brakóvannyy) — faulty, defective (of goods)}

\label{lesson:RU-C456-brakovannyy}



\begin{quote}
Before the new one: say the Russian for a foreign currency, then the Russian for a warranty.
\end{quote}

\subsection*{You'll want to know: \ru{бракованный}}
\textbf{\ru{бракованный}} (\emph{brakóvannyy}) — \textquotedblleft{}faulty, defective (of goods)\textquotedblright{}.

\begin{quote}
\ru{Мы} \ru{вернули} \ru{бракованный} \ru{чайник} \ru{в} \ru{магазин}.
\end{quote}

\begin{quote}
\emph{My vernúli brakóvannyy cháynik v magazín.}
\end{quote}

\begin{quote}
We returned the faulty kettle to the shop.
\end{quote}

Built on \emph{\ru{браковать}} — to reject as faulty.

\begin{cousinweb}
\emph{\ru{бракованный}} is built on \emph{\ru{браковать}} 'to reject as faulty', from \emph{\ru{брак}} 'a defect'. That \emph{\ru{брак}} is usually traced to Low German \emph{brak} 'damaged'; it is a different word from \emph{\ru{брак}} 'a marriage'. EtymDB lists \emph{\ru{браковать}} without a source.
\end{cousinweb}

\begin{grammarlens}[title={What you've built}]
One more word for telling, explaining and describing at B1.
\end{grammarlens}

\subsection*{Your turn}
\begin{itemize}
\raggedright
  \item \emph{Say it:} \emph{brakóvannyy}
  \item \emph{Say it:} \emph{brakóvannyy}, once more
  \item \emph{Say it:} a sentence of your own with \emph{brakóvannyy}
  \item \emph{Recall:} say the Russian for a foreign currency, then the Russian for a warranty, then say \emph{brakóvannyy} again
\end{itemize}

\begin{tcolorbox}[breakable,colback=teal!4,colframe=teal!35!black,title={Before you move on}]
What does \ru{бракованный} mean? (\textquotedblleft{}Faulty, defective (of goods)\textquotedblright{}.) Say it once more.
\end{tcolorbox}
\section[\texorpdfstring{\ru{экономить} (ekonómit)}{экономить (ekonómit)}]{\ru{экономить} (ekonómit) — to save (money or time), to economise}

\label{lesson:RU-C456-ekonomit}



\begin{quote}
Before the new one: say the Russian for a warranty, then the Russian for faulty.
\end{quote}

\subsection*{You'll want to know: \ru{экономить}}
\textbf{\ru{экономить}} (\emph{ekonómit}) — \textquotedblleft{}to save (money or time), to economise\textquotedblright{}.

\begin{quote}
\ru{Студенты} \ru{экономят} \ru{на} \ru{всём}.
\end{quote}

\begin{quote}
\emph{Studénty ekonómyat na vsyom.}
\end{quote}

\begin{quote}
Students save money on everything.
\end{quote}

Imperfective; the perfective is \emph{\ru{сэкономить}}. Present \emph{\ru{экономлю}, \ru{экономишь}}. Takes \emph{\ru{на}} + prepositional.

\begin{cousinweb}
\emph{\ru{экономить}} is built on \emph{\ru{экономия}} 'thrift', which is usually traced to Greek \emph{oikonomía} 'running a household'. EtymDB lists \emph{\ru{экономить}} and \emph{\ru{экономия}} without a source.
\end{cousinweb}

\begin{grammarlens}[title={What you've built}]
One more word for telling, explaining and describing at B1.
\end{grammarlens}

\subsection*{Your turn}
\begin{itemize}
\raggedright
  \item \emph{Say it:} \emph{ekonómit}
  \item \emph{Say it:} \emph{ekonómit}, once more
  \item \emph{Say it:} a sentence of your own with \emph{ekonómit}
  \item \emph{Recall:} say the Russian for a warranty, then the Russian for faulty, then say \emph{ekonómit} again
\end{itemize}

\begin{tcolorbox}[breakable,colback=teal!4,colframe=teal!35!black,title={Before you move on}]
What does \ru{экономить} mean? (\textquotedblleft{}To save (money or time), to economise\textquotedblright{}.) Say it once more.
\end{tcolorbox}
